%% kapa: exam=1 round=19 year=2008 period=1 form=A subject=accounting subject_ko=회계학 seq=32 number=112
%% source: Q-Net 감정평가사 기출문제 (cst00309) artlSeq=1001962 file=extracted/1008/1001962/2008년 제19회 감정평가사 1차시험 기출문제/1교시.hwp
%% conversion: kordoc 4.24.51 -> tools/gen1.py (auto); review=proofread-sonnet; flags=none
\begin{question}{112}
(주)하나는 두 사업부 X와 Y를 운영하고 있다. 사업부 X는 부품을 생산하고, 사업부 Y는 그 부품을 조립하여 완제품을 생산하고 있다. 부품과 완제품은 모두 외부시장이 존재하여 각 사업부는 이익중심점으로 운영되고 있고, 그 동안 부품의 대체가격은 장기평균시장가격으로 결정하였다. 관련자료는 다음과 같다.

{\small\centering
\begin{tabular}{lr}
\hline
완제품의 단위당 추정판매가격 & ￦220 \\
부품의 단위당 장기평균판매가격 & ￦150 \\
사업부 X의 단위당 변동원가 & ￦100 \\
사업부 Y의 단위당 추가가공원가 & ￦80 \\
\hline
\end{tabular}
\par}

현재 사업부 X에서는 다음의 두 대안을 고려하고 있다. \par

{\small\centering
\begin{tabular}{lp{0.72\linewidth}}
\hline
(대안1) & 외부판매가격을 단위당 ￦148으로 인하하고 10,000개의 부품을 시장판매한다. \\
(대안2) & 외부판매가격을 단위당 ￦150으로 하고 9,000개의 부품을 시장판매하며, 사업부 Y에 1,000개의 부품을 대체한다. \\
\hline
\end{tabular}
\par}

대안2에서 사업부 X는 대체가격을 얼마로 결정해야 대안1과 동일한 공헌이익을 얻을 수 있는가? \par
\begin{choices}
  \item ￦112
  \item ￦120
  \item ￦130
  \item ￦148
  \item ￦150
\end{choices}
\end{question}
\answer{3}
